\section{Analytics Methods}
The code below exists and is tested; the leads write the method justification.

\textbf{Detection (AN-1, AN-3).} EARS C2~\cite{hutwagner2003ears} runs on seasonally adjusted, completeness-scaled
weekly counts with Spark window functions; CUSUM runs per district and disease in \texttt{applyInPandas}.
Farrington~\cite{farrington1996} was ruled out because it needs about five years of history before its first
test week. Silent facility-weeks are treated as missing, not zero.

\textbf{Forecasting (AN-4).} One negative-binomial GLM per district for dengue and diarrhoea, with weather
lagged at least as far as the horizon, fitted in parallel through a cogrouped \texttt{applyInPandas}.

\textbf{Demand, allocation and calendar (AN-6, AN-7, AN-9).} Forecast quantiles become beds, IV saline, NS1 kits,
ORS and antimalarials with the assumptions in \texttt{analytics/forecast/demand\_assumptions.py}; the brief's LP
is solved with PuLP~3.3.2 and CBC; STL per division and disease gives the onset and peak weeks.

\begin{teammate}{Rafi}
Write the justification of each method from the literature, the LP formulation with our normalisation per
resource, and what $\alpha$ means for policy. Check the allocation first: in the SF=0.05 run the national net
need was 0 for beds, IV saline, NS1 kits and ORS, so the LP had almost nothing to allocate.
\end{teammate}
